\chapter{Data landscape and pipelines}
\label{ch:data-landscape}

\begin{quote}
\emph{``Data is the new oil --- but like oil, it must be refined before it is useful.''}
--- Clive Humby
\end{quote}

% ============================================================
\section{What is data preprocessing?}

Data preprocessing is the set of transformations applied to raw data before analysis or modeling. It bridges the gap between the data you \emph{have} and the data you \emph{need}.

A typical preprocessing pipeline includes:
\begin{enumerate}
  \item \textbf{Loading} --- reading data from files, databases, or APIs.
  \item \textbf{Inspection} --- understanding structure, types, and quality.
  \item \textbf{Cleaning} --- handling missing values, duplicates, and errors.
  \item \textbf{Transformation} --- encoding, scaling, and creating features.
  \item \textbf{Validation} --- confirming the output meets downstream requirements.
\end{enumerate}

\begin{datatip}
In industry surveys, data practitioners report spending 60--80\% of project time on preprocessing. Mastering this phase is the single highest-leverage skill in data science.
\end{datatip}

% ============================================================
\section{Data types and structures}

\subsection{Structured, semi-structured, and unstructured data}

\begin{itemize}
  \item \textbf{Structured:} tables with rows and columns (CSV, SQL databases, Excel).
  \item \textbf{Semi-structured:} data with flexible schemas (JSON, XML, YAML).
  \item \textbf{Unstructured:} free text, images, audio, video.
\end{itemize}

This course focuses on structured and semi-structured data, with one chapter (Chapter~7) devoted to text.

\subsection{Variable types}

\begin{minted}{python}
import pandas as pd

# Load a real dataset to inspect variable types
url = ("https://raw.githubusercontent.com/datasciencedojo/"
       "datasets/master/titanic.csv")
df = pd.read_csv(url)
print(df.dtypes)
\end{minted}

Output (abbreviated):
\begin{minted}{text}
PassengerId      int64
Survived         int64
Pclass           int64
Name            object
Sex             object
Age            float64
SibSp            int64
Fare           float64
Embarked        object
\end{minted}

\begin{preprocesstip}
The \texttt{object} dtype in pandas usually means string data, but it can also hide mixed types. Always verify with \texttt{df[col].apply(type).value\_counts()}.
\end{preprocesstip}

\subsection{Numerical vs.\ categorical}

\begin{minted}{python}
numerical = df.select_dtypes(include="number").columns.tolist()
categorical = df.select_dtypes(include="object").columns.tolist()

print(f"Numerical columns ({len(numerical)}): {numerical}")
print(f"Categorical columns ({len(categorical)}): {categorical}")
\end{minted}

% ============================================================
\section{Common data formats}

\begin{longtable}{lllp{5cm}}
\toprule
\textbf{Format} & \textbf{Extension} & \textbf{pandas reader} & \textbf{Notes} \\
\midrule
CSV & \texttt{.csv} & \texttt{read\_csv} & Most common; watch for delimiter and encoding \\
Excel & \texttt{.xlsx} & \texttt{read\_excel} & Requires \texttt{openpyxl}; multi-sheet support \\
JSON & \texttt{.json} & \texttt{read\_json} & Nested structures need \texttt{json\_normalize} \\
Parquet & \texttt{.parquet} & \texttt{read\_parquet} & Columnar, fast, type-preserving \\
SQL & --- & \texttt{read\_sql} & Requires \texttt{sqlalchemy} connection \\
\bottomrule
\end{longtable}

% ============================================================
\section{First look at a real dataset}

We will use the Melbourne Housing dataset throughout several chapters:

\begin{minted}{python}
# Melbourne Housing dataset
melb = pd.read_csv("melb_data.csv")

# Shape: rows x columns
print(f"Shape: {melb.shape}")

# First 3 rows
print(melb.head(3))

# Summary statistics
print(melb.describe())

# Missing values per column
print(melb.isnull().sum().sort_values(ascending=False).head(10))
\end{minted}

\begin{minted}{python}
# Quick quality report
def data_quality_report(df):
    """Generate a quick quality report for any DataFrame."""
    report = pd.DataFrame({
        "dtype": df.dtypes,
        "non_null": df.notnull().sum(),
        "null_count": df.isnull().sum(),
        "null_pct": (df.isnull().sum() / len(df) * 100).round(1),
        "n_unique": df.nunique(),
        "sample_value": df.iloc[0]
    })
    return report.sort_values("null_pct", ascending=False)

print(data_quality_report(melb))
\end{minted}

\begin{datatip}
Always run a quality report \emph{before} any analysis. The five numbers you need immediately are: shape, dtypes, null counts, unique counts, and summary statistics.
\end{datatip}

% ============================================================
\section{The preprocessing workflow}

\begin{minted}{python}
# A minimal end-to-end preprocessing workflow
import pandas as pd
import numpy as np

# 1. Load
df = pd.read_csv("melb_data.csv")

# 2. Inspect
print(df.info())

# 3. Clean: drop duplicates, handle missing values
df = df.drop_duplicates()
df["Car"] = df["Car"].fillna(df["Car"].median())
df = df.dropna(subset=["Price"])

# 4. Transform: encode categorical, scale numerical
df["Type_encoded"] = df["Type"].map({"h": 0, "u": 1, "t": 2})

# 5. Validate
assert df.isnull().sum().sum() == 0 or True  # relaxed for demo
print(f"Clean shape: {df.shape}")
\end{minted}

% ============================================================
\section{Pandas review: essential operations}

\subsection{Selection and filtering}

\begin{minted}{python}
# Select columns
prices = melb[["Suburb", "Price", "Rooms"]]

# Filter rows
expensive = melb[melb["Price"] > 2_000_000]
south = melb[melb["Regionname"] == "Southern Metropolitan"]

# Combined
big_south = melb[(melb["Regionname"] == "Southern Metropolitan") &
                 (melb["Rooms"] >= 4)]
print(f"Large southern houses: {len(big_south)}")
\end{minted}

\subsection{Groupby and aggregation}

\begin{minted}{python}
# Average price by region
region_stats = (melb.groupby("Regionname")["Price"]
                .agg(["mean", "median", "count"])
                .sort_values("median", ascending=False))
print(region_stats)
\end{minted}

\begin{warningbox}
Pandas \texttt{groupby} silently drops \texttt{NaN} group keys by default. If your grouping column has missing values, they will be excluded from the result. Use \texttt{dropna=False} to include them.
\end{warningbox}

% ============================================================
\section{Tools and libraries overview}

\begin{minted}{python}
# Core preprocessing stack
import pandas as pd               # data manipulation
import numpy as np                 # numerical operations
import matplotlib.pyplot as plt    # visualization
import seaborn as sns              # statistical plots
import missingno as msno           # missing data visualization
from sklearn.preprocessing import StandardScaler, OneHotEncoder
from sklearn.impute import SimpleImputer, KNNImputer
from sklearn.pipeline import Pipeline
from sklearn.compose import ColumnTransformer
\end{minted}

% ============================================================
\section{Exercises}

\begin{exercise}
\begin{enumerate}
  \item Download the Melbourne Housing dataset. Write a quality report showing: number of rows, columns, dtype of each column, percentage of missing values, and number of unique values.
  \item Load the Gapminder dataset (\texttt{gapminder.csv} or via \texttt{plotly.express}). Identify which columns are numerical, which are categorical, and which have missing values.
  \item Write a function \texttt{detect\_mixed\_types(df)} that checks each column for mixed Python types and returns a list of problematic columns.
  \item Compare loading the Titanic dataset as CSV vs.\ Parquet. Measure file size and load time using \texttt{\%timeit}.
\end{enumerate}
\end{exercise}

% ============================================================
\section{Chapter summary}

\begin{itemize}
  \item Data preprocessing transforms raw data into analysis-ready data and consumes the majority of project time.
  \item Data comes in structured, semi-structured, and unstructured formats; pandas handles the first two.
  \item Every preprocessing workflow follows five steps: load, inspect, clean, transform, validate.
  \item A quick quality report (shape, dtypes, nulls, uniques, statistics) should be your first action on any new dataset.
  \item The core Python stack for preprocessing is pandas, numpy, scikit-learn, matplotlib, and missingno.
\end{itemize}
